% Worked Examples: Topic 3.9 - Robust Performance & Fundamental Limits
\documentclass[11pt,a4paper]{article}
\usepackage{../worked-examples-template}

\title{\textbf{Worked Examples}\\[0.5em]
\large Topic 3.9: Robust Performance \& Fundamental Limits\\[0.3em]
\normalsize \coursesubtitle{} -- \coursetitle}
\author{Assoc.\ Prof.\ William Robertson \and Dr Sean McGowan}
\date{\today}

\begin{document}

\maketitle
\tableofcontents
\newpage

\section{Concept 3.9.1: Uncertainty Modelling}

\subsection{Example 1: Parametric Uncertainty in Cruise Control}

Understanding parametric uncertainty is crucial for designing robust control systems. This example demonstrates how parameter variations affect controller performance.

\paragraph{Problem:} 
A cruise control system is designed with a PI controller for a car with nominal mass $m_0 = 1600$ kg. The linearised vehicle dynamics around a constant speed are given by:
\[
m\ddot{x} = F - c\dot{x} - mg\sin\theta
\]
where $x$ is position, $F$ is the drive force, $c = 10$ N·s/m is the drag coefficient, $g = 9.81$ m/s$^2$, and $\theta$ is the road slope.

The PI controller is designed as $C(s) = k_p + k_i/s$ with $k_p = 50$ and $k_i = 20$ for the nominal mass. 

(a) Determine the closed-loop characteristic equation for arbitrary mass $m$.

(b) Calculate the maximum mass $m_{\max}$ for which the closed-loop system remains stable.

(c) Explain why parametric uncertainty matters for robust control design.

\paragraph{Solution:}

\textbf{Step 1: System transfer function}

For velocity control at constant speed, let $v = \dot{x}$ and linearise around a slope angle $\theta$. The transfer function from force to velocity is:
\[
P(s) = \frac{V(s)}{F(s)} = \frac{1}{ms + c}
\]

For $m_0 = 1600$ kg and $c = 10$:
\[
P_0(s) = \frac{1}{1600s + 10} = \frac{1/1600}{s + 0.00625}
\]

\textbf{Step 2: Closed-loop characteristic equation}

With PI controller $C(s) = k_p + k_i/s = (k_ps + k_i)/s$, the loop transfer function is:
\[
L(s) = P(s)C(s) = \frac{1}{ms + c} \cdot \frac{k_ps + k_i}{s} = \frac{k_ps + k_i}{s(ms + c)}
\]

The closed-loop characteristic equation is $1 + L(s) = 0$:
\[
s(ms + c) + k_ps + k_i = 0
\]
\[
ms^2 + (c + k_p)s + k_i = 0
\]

\textbf{Step 3: Stability analysis}

For stability, all coefficients must be positive (necessary condition for 2nd order system):
\begin{itemize}
\item $m > 0$ (always satisfied physically)
\item $c + k_p > 0 \Rightarrow k_p > -c$ (satisfied since $k_p = 50 > 0$)
\item $k_i > 0$ (satisfied since $k_i = 20 > 0$)
\end{itemize}

For a second-order system, if all coefficients are positive, the system is stable. Therefore, the system remains stable for all positive masses.

\textbf{Step 4: Performance degradation}

While the system remains stable, performance degrades with increasing mass. The natural frequency and damping ratio are:
\[
\omega_n = \sqrt{\frac{k_i}{m}}, \quad \zeta = \frac{c + k_p}{2\sqrt{mk_i}}
\]

For $m = 1600$ kg (nominal):
\[
\omega_n = \sqrt{\frac{20}{1600}} = 0.112 \text{ rad/s}, \quad \zeta = \frac{10 + 50}{2\sqrt{1600 \cdot 20}} = \frac{60}{253} = 0.237
\]

For $m = 2000$ kg (25\% increase):
\[
\omega_n = \sqrt{\frac{20}{2000}} = 0.100 \text{ rad/s}, \quad \zeta = \frac{60}{2\sqrt{2000 \cdot 20}} = \frac{60}{283} = 0.212
\]

The response becomes slower (lower $\omega_n$) and slightly less damped (lower $\zeta$), leading to longer settling time and potentially larger overshoot.

\textbf{Key insights:}
\begin{itemize}
\item Parametric uncertainty doesn't necessarily cause instability but degrades performance
\item Controllers designed for nominal conditions may perform poorly under parameter variations
\item Robust design requires considering the full range of expected parameter values
\end{itemize}

\subsection{Example 2: Describing Unmodeled Dynamics with Multiplicative Uncertainty}

Unmodeled dynamics represent system behaviors not captured by our nominal model. This example shows how to characterize them mathematically.

\paragraph{Problem:}
A nominal process model is $P(s) = 1/(s+1)$. However, the actual system has unmodeled high-frequency dynamics and can be represented as $P_{\text{actual}}(s) = 1/[(s+1)(0.1s+1)]$.

(a) Express the actual system using multiplicative uncertainty: $P_{\text{actual}} = P(1 + \delta)$.

(b) Find the multiplicative uncertainty $\delta(s)$ and plot $|\delta(\ii\omega)|$.

(c) Determine the maximum value of $|\delta(\ii\omega)|$ and the frequency at which it occurs.

\paragraph{Solution:}

\textbf{Step 1: Calculate the multiplicative perturbation}

By definition of multiplicative uncertainty:
\[
P_{\text{actual}}(s) = P(s)(1 + \delta(s))
\]

Solving for $\delta(s)$:
\[
\delta(s) = \frac{P_{\text{actual}}(s)}{P(s)} - 1 = \frac{P_{\text{actual}}(s) - P(s)}{P(s)}
\]

Substituting the given models:
\[
\delta(s) = \frac{\frac{1}{(s+1)(0.1s+1)} - \frac{1}{s+1}}{\frac{1}{s+1}} = \frac{1}{(s+1)(0.1s+1)} \cdot (s+1) - 1
\]
\[
= \frac{1}{0.1s+1} - 1 = \frac{1 - (0.1s+1)}{0.1s+1} = \frac{-0.1s}{0.1s+1}
\]

\textbf{Step 2: Frequency response of the uncertainty}

At frequency $\omega$, substitute $s = \ii\omega$:
\[
\delta(\ii\omega) = \frac{-0.1\ii\omega}{0.1\ii\omega + 1} = \frac{-\ii 0.1\omega}{1 + \ii 0.1\omega}
\]

The magnitude is:
\[
|\delta(\ii\omega)| = \left|\frac{-\ii 0.1\omega}{1 + \ii 0.1\omega}\right| = \frac{0.1\omega}{\sqrt{1 + (0.1\omega)^2}}
\]

\textbf{Step 3: Analyze the uncertainty bound}

To find the maximum, take the derivative with respect to $\omega$ and set to zero. Alternatively, observe that:
\[
|\delta(\ii\omega)| = \frac{0.1\omega}{\sqrt{1 + 0.01\omega^2}}
\]

As $\omega \to 0$: $|\delta| \to 0$

As $\omega \to \infty$: $|\delta| \to \frac{0.1\omega}{0.1\omega} = 1$

The function is monotonically increasing and approaches 1 asymptotically.

At $\omega = 10$ rad/s:
\[
|\delta(\ii \cdot 10)| = \frac{1}{\sqrt{1 + 1}} = \frac{1}{\sqrt{2}} = 0.707
\]

At $\omega = 100$ rad/s:
\[
|\delta(\ii \cdot 100)| = \frac{10}{\sqrt{1 + 100}} = \frac{10}{\sqrt{101}} = 0.995
\]

\textbf{Step 4: Physical interpretation}

The multiplicative uncertainty $\delta(s) = -0.1s/(0.1s+1)$ represents:
\begin{itemize}
\item A high-pass characteristic (magnitude increases with frequency)
\item Maximum relative error approaching 100\% at very high frequencies
\item The unmodeled pole at $s = -10$ causes phase lag and gain reduction at high frequencies
\end{itemize}

\textbf{Key insights:}
\begin{itemize}
\item Multiplicative uncertainty naturally increases with frequency for unmodeled fast dynamics
\item The uncertainty bound $|\delta(\ii\omega)| < 1$ for all $\omega$ means the actual plant is within 100\% of the nominal at high frequencies
\item For robust stability, the complementary sensitivity $|T(\ii\omega)|$ must be small where $|\delta(\ii\omega)|$ is large
\end{itemize}

\section{Concept 3.9.2: Robust Stability}

\subsection{Example 1: Stability Bounds with Multiplicative Perturbations}

Robust stability ensures a control system remains stable despite model uncertainties. This example applies the multiplicative uncertainty condition.

\paragraph{Problem:}
Consider a feedback system with nominal process $P(s) = 5/(s+2)$ and controller $C(s) = 2$. 

(a) Calculate the complementary sensitivity function $T(s)$.

(b) Determine the maximum magnitude of $T$ over all frequencies: $M_t = \sup_\omega |T(\ii\omega)|$.

(c) Using the robust stability condition $|\delta(\ii\omega)| < 1/|T(\ii\omega)|$, find the maximum allowable multiplicative uncertainty at $\omega = 0$, $\omega = 2$, and $\omega = 10$ rad/s.

(d) Interpret the results for controller design.

\paragraph{Solution:}

\textbf{Step 1: Loop transfer function and complementary sensitivity}

The loop transfer function is:
\[
L(s) = P(s)C(s) = \frac{5}{s+2} \cdot 2 = \frac{10}{s+2}
\]

The complementary sensitivity function is:
\[
T(s) = \frac{L(s)}{1 + L(s)} = \frac{10/(s+2)}{1 + 10/(s+2)} = \frac{10}{s+2+10} = \frac{10}{s+12}
\]

\textbf{Step 2: Frequency response of $T$}

At frequency $\omega$, substitute $s = \ii\omega$:
\[
T(\ii\omega) = \frac{10}{\ii\omega + 12}
\]

The magnitude is:
\[
|T(\ii\omega)| = \frac{10}{\sqrt{\omega^2 + 144}}
\]

\textbf{Step 3: Find maximum magnitude $M_t$}

Since $|T(\ii\omega)|$ is monotonically decreasing with $\omega$:
\[
M_t = \max_\omega |T(\ii\omega)| = |T(\ii \cdot 0)| = \frac{10}{12} = 0.833
\]

This occurs at $\omega = 0$ (DC gain).

\textbf{Step 4: Robust stability bounds at specific frequencies}

The robust stability condition states: $|\delta(\ii\omega)| < 1/|T(\ii\omega)|$

At $\omega = 0$ rad/s:
\[
|T(\ii \cdot 0)| = \frac{10}{12} = 0.833 \quad \Rightarrow \quad |\delta| < \frac{1}{0.833} = 1.20
\]

At $\omega = 2$ rad/s:
\[
|T(\ii \cdot 2)| = \frac{10}{\sqrt{4 + 144}} = \frac{10}{\sqrt{148}} = 0.822 \quad \Rightarrow \quad |\delta| < 1.22
\]

At $\omega = 10$ rad/s:
\[
|T(\ii \cdot 10)| = \frac{10}{\sqrt{100 + 144}} = \frac{10}{\sqrt{244}} = 0.640 \quad \Rightarrow \quad |\delta| < 1.56
\]

\textbf{Step 5: Interpretation}

\begin{itemize}
\item At low frequencies ($\omega \approx 0$), the uncertainty bound is tighter: $|\delta| < 1.20$
\item At high frequencies ($\omega = 10$), larger uncertainties are tolerable: $|\delta| < 1.56$
\item This makes sense: the controller has more gain at low frequencies, making the system more sensitive to model errors
\item The maximum uncertainty across all frequencies is limited by the minimum of $1/|T(\ii\omega)|$, which equals $1/M_t = 1.20$
\end{itemize}

\textbf{Key insights:}
\begin{itemize}
\item Keeping $M_t$ small (equivalently, the sensitivity peak $M_s$ small) improves robustness
\item The complementary sensitivity $T$ determines tolerance to high-frequency uncertainty
\item Trade-off: high controller gain improves disturbance rejection but reduces robustness margins
\end{itemize}

\subsection{Example 2: Performance Under Uncertainty - Load Disturbance Attenuation}

Understanding how uncertainty affects disturbance rejection is essential for robust performance analysis.

\paragraph{Problem:}
A process has nominal model $P_0(s) = 1/(s+1)$ with multiplicative uncertainty $\delta$ such that the actual process is $P(s) = P_0(s)(1+\delta(s))$. A controller $C(s) = 5$ is used.

(a) Calculate the sensitivity function $S(s)$ for the nominal system.

(b) Determine the transfer function $G_{yv}(s)$ from load disturbance $v$ to output $y$ for the nominal system.

(c) If the actual process has $\delta(s) = 0.2$, find the actual $G_{yv}$ and compare with the nominal case.

(d) Calculate the relative error in disturbance rejection performance.

\paragraph{Solution:}

\textbf{Step 1: Nominal sensitivity function}

With nominal process $P_0(s) = 1/(s+1)$ and controller $C(s) = 5$:
\[
L_0(s) = P_0(s)C(s) = \frac{5}{s+1}
\]

The sensitivity function is:
\[
S_0(s) = \frac{1}{1 + L_0(s)} = \frac{1}{1 + 5/(s+1)} = \frac{s+1}{s+1+5} = \frac{s+1}{s+6}
\]

\textbf{Step 2: Load disturbance transfer function (nominal)}

The transfer function from load disturbance to output is:
\[
G_{yv,0}(s) = P_0(s)S_0(s) = \frac{1}{s+1} \cdot \frac{s+1}{s+6} = \frac{1}{s+6}
\]

\textbf{Step 3: Actual system with uncertainty}

With $\delta = 0.2$, the actual process is:
\[
P(s) = P_0(s)(1 + \delta) = \frac{1}{s+1} \cdot 1.2 = \frac{1.2}{s+1}
\]

The actual loop transfer function:
\[
L(s) = P(s)C(s) = \frac{1.2}{s+1} \cdot 5 = \frac{6}{s+1}
\]

The actual sensitivity:
\[
S(s) = \frac{1}{1 + 6/(s+1)} = \frac{s+1}{s+7}
\]

The actual disturbance transfer function:
\[
G_{yv}(s) = P(s)S(s) = \frac{1.2}{s+1} \cdot \frac{s+1}{s+7} = \frac{1.2}{s+7}
\]

\textbf{Step 4: Relative error analysis}

The relative change in $G_{yv}$ can be computed. At steady state ($s \to 0$):

Nominal: $G_{yv,0}(0) = 1/6 = 0.167$

Actual: $G_{yv}(0) = 1.2/7 = 0.171$

Relative error:
\[
\frac{G_{yv}(0) - G_{yv,0}(0)}{G_{yv,0}(0)} = \frac{0.171 - 0.167}{0.167} = 0.024 = 2.4\%
\]

Using the theoretical relation $\frac{\mathrm{d}G_{yv}}{G_{yv}} = S\frac{\mathrm{d}P}{P}$:

At low frequency where $S_0(0) = 1/6$:
\[
\frac{\Delta G_{yv}}{G_{yv}} \approx S_0(0) \cdot \frac{\Delta P}{P} = \frac{1}{6} \cdot 0.2 = 0.033 = 3.3\%
\]

The small discrepancy is due to the linearization approximation.

\textbf{Key insights:}
\begin{itemize}
\item Load disturbance attenuation is relatively insensitive to process variations at frequencies where $S$ is small
\item The sensitivity function $S$ acts as a "scaling factor" for how process uncertainty affects disturbance rejection
\item Good feedback design (small $S$ at low frequencies) provides both good disturbance rejection AND insensitivity to model errors
\end{itemize}

\section{Concept 3.9.3: System Design}

\subsection{Example 1: Stabilisability Check Using the Popov-Belevitch-Hautus Test}

The PBH test provides an algebraic method to determine if a system can be stabilised. This is fundamental for determining feasibility of control design.

\paragraph{Problem:}
Consider a system with state-space representation:
\[
A = \begin{bmatrix} 2 & 1 & 0 \\ 0 & -1 & 1 \\ 0 & 0 & 3 \end{bmatrix}, \quad
B = \begin{bmatrix} 0 \\ 1 \\ 0 \end{bmatrix}
\]

(a) Determine if the system is stabilisable using the Popov-Belevitch-Hautus (PBH) test.

(b) Identify which eigenvalues can be moved by state feedback.

(c) Explain why this system can or cannot be stabilised.

\paragraph{Solution:}

\textbf{Step 1: Find eigenvalues of $A$}

The eigenvalues are found from $\det(A - \lambda I) = 0$. Since $A$ is upper triangular, the eigenvalues are the diagonal elements:
\[
\lambda_1 = 2, \quad \lambda_2 = -1, \quad \lambda_3 = 3
\]

Eigenvalues in the right half-plane (unstable): $\lambda_1 = 2$ and $\lambda_3 = 3$.

\textbf{Step 2: Apply PBH test for stabilisability}

The PBH test states that a system is stabilisable if and only if:
\[
\operatorname{rank}([A - \lambda I \quad B]) = n
\]
for all eigenvalues $\lambda$ with $\operatorname{Re}(\lambda) \geq 0$ (right half-plane).

We need to check this for $\lambda_1 = 2$ and $\lambda_3 = 3$.

\textbf{For $\lambda_1 = 2$:}
\[
A - 2I = \begin{bmatrix} 0 & 1 & 0 \\ 0 & -3 & 1 \\ 0 & 0 & 1 \end{bmatrix}
\]
\[
[A - 2I \quad B] = \begin{bmatrix} 0 & 1 & 0 & 0 \\ 0 & -3 & 1 & 1 \\ 0 & 0 & 1 & 0 \end{bmatrix}
\]

Row reduce to find rank:
\[
\begin{bmatrix} 0 & 1 & 0 & 0 \\ 0 & -3 & 1 & 1 \\ 0 & 0 & 1 & 0 \end{bmatrix} \sim \begin{bmatrix} 0 & 1 & 0 & 0 \\ 0 & 0 & 1 & 1 \\ 0 & 0 & 1 & 0 \end{bmatrix} \sim \begin{bmatrix} 0 & 1 & 0 & 0 \\ 0 & 0 & 1 & 0 \\ 0 & 0 & 0 & 1 \end{bmatrix}
\]

Rank = 3 ✓ (satisfies PBH test for $\lambda = 2$)

\textbf{For $\lambda_3 = 3$:}
\[
A - 3I = \begin{bmatrix} -1 & 1 & 0 \\ 0 & -4 & 1 \\ 0 & 0 & 0 \end{bmatrix}
\]
\[
[A - 3I \quad B] = \begin{bmatrix} -1 & 1 & 0 & 0 \\ 0 & -4 & 1 & 1 \\ 0 & 0 & 0 & 0 \end{bmatrix}
\]

Row reduce:
\[
\begin{bmatrix} -1 & 1 & 0 & 0 \\ 0 & -4 & 1 & 1 \\ 0 & 0 & 0 & 0 \end{bmatrix}
\]

Rank = 2 ✗ (fails PBH test for $\lambda = 3$)

\textbf{Step 3: Conclusion}

The system is \textbf{NOT stabilisable} because the PBH test fails for the unstable eigenvalue $\lambda_3 = 3$.

Physically, this means:
\begin{itemize}
\item The eigenvalue $\lambda = 2$ is reachable and can be moved by state feedback
\item The eigenvalue $\lambda = 3$ is \emph{unreachable} from the input and cannot be affected by any control
\item Since $\lambda = 3$ is unstable and cannot be controlled, the system cannot be stabilised
\end{itemize}

\textbf{Key insights:}
\begin{itemize}
\item Stabilisability requires that all unstable modes be reachable from the input
\item A system can have unreachable modes, but only if they are stable
\item The PBH test provides a simple algebraic check without computing eigenvectors
\end{itemize}

\subsection{Example 2: Bode's Integral Formula and the Waterbed Effect}

Bode's Integral Formula reveals fundamental limitations: reducing sensitivity at one frequency inevitably increases it elsewhere.

\paragraph{Problem:}
Consider a feedback system with an unstable process $P(s) = 1/(s-1)$ and proportional controller $C(s) = k$. 

(a) For $k = 5$, calculate the sensitivity function $S(s)$ and loop transfer function $L(s)$.

(b) Evaluate $\log|S(\ii\omega)|$ at $\omega = 0.5$, $\omega = 2$, and $\omega = 10$ rad/s.

(c) Apply Bode's Integral Formula to find $\int_0^\infty \log|S(\ii\omega)|\,\mathrm{d}\omega$.

(d) Explain the waterbed effect for this system.

\paragraph{Solution:}

\textbf{Step 1: Loop transfer function and sensitivity}

Loop transfer function:
\[
L(s) = P(s)C(s) = \frac{k}{s-1} = \frac{5}{s-1}
\]

Sensitivity function:
\[
S(s) = \frac{1}{1 + L(s)} = \frac{1}{1 + 5/(s-1)} = \frac{s-1}{s-1+5} = \frac{s-1}{s+4}
\]

\textbf{Step 2: Evaluate $\log|S(\ii\omega)|$ at specific frequencies}

At $\omega = 0.5$ rad/s:
\[
S(\ii \cdot 0.5) = \frac{\ii \cdot 0.5 - 1}{\ii \cdot 0.5 + 4} = \frac{-1 + \ii \cdot 0.5}{4 + \ii \cdot 0.5}
\]
\[
|S(\ii \cdot 0.5)| = \frac{\sqrt{1 + 0.25}}{\sqrt{16 + 0.25}} = \frac{\sqrt{1.25}}{\sqrt{16.25}} = \frac{1.118}{4.031} = 0.277
\]
\[
\log|S(\ii \cdot 0.5)| = \log(0.277) = -1.284
\]

At $\omega = 2$ rad/s:
\[
|S(\ii \cdot 2)| = \frac{\sqrt{1 + 4}}{\sqrt{16 + 4}} = \frac{\sqrt{5}}{\sqrt{20}} = \frac{2.236}{4.472} = 0.500
\]
\[
\log|S(\ii \cdot 2)| = \log(0.5) = -0.693
\]

At $\omega = 10$ rad/s:
\[
|S(\ii \cdot 10)| = \frac{\sqrt{1 + 100}}{\sqrt{16 + 100}} = \frac{\sqrt{101}}{\sqrt{116}} = \frac{10.05}{10.77} = 0.933
\]
\[
\log|S(\ii \cdot 10)| = \log(0.933) = -0.069
\]

\textbf{Step 3: Apply Bode's Integral Formula}

Bode's Integral Formula states:
\[
\int_0^\infty \log|S(\ii\omega)|\,\mathrm{d}\omega = \pi \sum p_k
\]
where $p_k$ are the right half-plane poles of $L(s)$.

For our system, $L(s) = 5/(s-1)$ has one right half-plane pole at $p_1 = 1$.

Therefore:
\[
\int_0^\infty \log|S(\ii\omega)|\,\mathrm{d}\omega = \pi \cdot 1 = \pi \approx 3.14
\]

\textbf{Step 4: Interpretation - The waterbed effect}

\begin{itemize}
\item The integral is \emph{constant} and positive for unstable systems
\item At low frequencies ($\omega = 0.5$), $\log|S| < 0$ means $|S| < 1$ (good disturbance rejection)
\item At high frequencies ($\omega = 10$), $\log|S| > 0$ means $|S| > 1$ (amplification)
\item The negative area (good performance) at low frequencies must be balanced by positive area (poor performance) at higher frequencies
\item This is the "waterbed effect": pushing down sensitivity at one frequency causes it to pop up elsewhere
\end{itemize}

If we tried to increase $k$ to improve low-frequency disturbance rejection, we would:
\begin{itemize}
\item Reduce $|S|$ at low frequencies (more negative contribution to integral)
\item Force $|S|$ to increase more at mid-range frequencies (to maintain constant integral)
\item Create a larger sensitivity peak, reducing robustness
\end{itemize}

\textbf{Key insights:}
\begin{itemize}
\item Unstable systems have inherent performance limitations
\item The total "cost" of disturbance rejection is fixed by the unstable pole location
\item Control design is always a compromise between performance at different frequencies
\item Stable systems ($\sum p_k = 0$) have zero integral, allowing better performance flexibility
\end{itemize}

\section{Concept 3.9.4: Robust Control}

\subsection{Example 1: PI Controller Design with Fast Stable Pole}

When a process has stable poles faster than the desired closed-loop bandwidth, special care is needed to avoid robustness problems.

\paragraph{Problem:}
Consider a first-order process $P(s) = 2/(s+10)$ with a fast stable pole at $s = -10$. Design a PI controller $C(s) = k_p + k_i/s$ to place the closed-loop poles at $s = -2 \pm \ii$.

(a) Calculate the required controller gains $k_p$ and $k_i$.

(b) Determine the sensitivity function $S(s)$ and complementary sensitivity function $T(s)$.

(c) Find the maximum sensitivity $M_s = \max_\omega |S(\ii\omega)|$ and maximum complementary sensitivity $M_t = \max_\omega |T(\ii\omega)|$.

(d) Explain why this design has poor robustness and suggest an improvement.

\paragraph{Solution:}

\textbf{Step 1: Closed-loop characteristic equation}

The process transfer function is $P(s) = 2/(s+10)$ (where $b = 2$ and $a = 10$).

The PI controller is $C(s) = k_p + k_i/s = (k_ps + k_i)/s$.

The loop transfer function:
\[
L(s) = P(s)C(s) = \frac{2}{s+10} \cdot \frac{k_ps + k_i}{s} = \frac{2k_ps + 2k_i}{s(s+10)}
\]

The closed-loop characteristic equation is:
\[
1 + L(s) = 0 \quad \Rightarrow \quad s(s+10) + 2k_ps + 2k_i = 0
\]
\[
s^2 + (10 + 2k_p)s + 2k_i = 0
\]

\textbf{Step 2: Pole placement calculation}

For poles at $s = -2 \pm \ii$, the desired characteristic polynomial is:
\[
(s - (-2 + \ii))(s - (-2 - \ii)) = (s+2-\ii)(s+2+\ii) = (s+2)^2 + 1
\]
\[
= s^2 + 4s + 4 + 1 = s^2 + 4s + 5
\]

Matching coefficients with $s^2 + (10 + 2k_p)s + 2k_i = 0$:
\[
10 + 2k_p = 4 \quad \Rightarrow \quad k_p = -3
\]
\[
2k_i = 5 \quad \Rightarrow \quad k_i = 2.5
\]

\textbf{Step 3: Sensitivity functions}

The controller is $C(s) = -3 + 2.5/s = (-3s + 2.5)/s$.

Notice that $k_p < 0$, which means the controller has a zero in the right half-plane:
\[
C(s) = \frac{-3s + 2.5}{s} = -3\frac{s - 0.833}{s}
\]

The zero is at $s = 2.5/3 = 0.833$ (RHP zero).

Loop transfer function:
\[
L(s) = \frac{2(-3s + 2.5)}{s(s+10)} = \frac{-6s + 5}{s(s+10)}
\]

Sensitivity function:
\[
S(s) = \frac{1}{1 + L(s)} = \frac{s(s+10)}{s(s+10) + (-6s + 5)} = \frac{s(s+10)}{s^2 + 4s + 5}
\]

Complementary sensitivity:
\[
T(s) = \frac{L(s)}{1 + L(s)} = \frac{-6s + 5}{s^2 + 4s + 5}
\]

\textbf{Step 4: Calculate maximum sensitivities}

For $M_s$, we need to find $\max_\omega |S(\ii\omega)|$:
\[
S(\ii\omega) = \frac{\ii\omega(\ii\omega + 10)}{(\ii\omega)^2 + 4\ii\omega + 5} = \frac{-\omega^2 + 10\ii\omega}{-\omega^2 + 4\ii\omega + 5}
\]
\[
= \frac{(5-\omega^2) + 10\ii\omega}{(5-\omega^2) + 4\ii\omega}
\]

The magnitude is:
\[
|S(\ii\omega)| = \frac{\sqrt{(5-\omega^2)^2 + 100\omega^2}}{\sqrt{(5-\omega^2)^2 + 16\omega^2}}
\]

At the resonant frequency $\omega_r = \sqrt{5} \approx 2.24$ rad/s (where the denominator is minimized):
\[
|S(\ii\sqrt{5})| = \frac{\sqrt{0 + 500}}{\sqrt{0 + 80}} = \frac{\sqrt{500}}{\sqrt{80}} = \frac{22.36}{8.94} = 2.50
\]

Therefore, $M_s \approx 2.50$ (very large, indicating poor robustness!).

\textbf{Step 5: Analysis and improvement}

\textbf{Why is robustness poor?}
\begin{itemize}
\item The fast stable pole at $s = -10$ is much faster than the desired bandwidth ($\omega_n = \sqrt{5} \approx 2.24$)
\item To place slow poles at $s = -2 \pm \ii$, we need $k_p < 0$, creating an RHP zero in the controller
\item This RHP zero causes a large sensitivity peak: $M_s = 2.50 > 2$, indicating poor stability margins
\item The minimum robust stability margin is $1/M_s = 0.4$, meaning only 40\% multiplicative uncertainty is tolerable
\end{itemize}

\textbf{Improvement strategy:}
\begin{itemize}
\item Choose closed-loop poles faster than the process pole: e.g., $s = -15 \pm 5\ii$
\item This ensures $k_p > 0$ and avoids the RHP zero
\item Accept the natural fast dynamics of the process rather than fighting them
\item Alternatively, use a lead compensator to cancel the fast pole (if exact cancellation is feasible)
\end{itemize}

\textbf{Key insights:}
\begin{itemize}
\item Attempting to make the closed-loop slower than a fast stable process pole leads to poor robustness
\item Negative proportional gain creates RHP zeros that amplify sensitivity
\item Robust design: let fast stable poles remain fast in closed-loop
\end{itemize}

\subsection{Example 2: Pole Placement with Slow Stable Zero}

Process zeros slower than the desired bandwidth also require careful treatment to maintain robustness.

\paragraph{Problem:}
Consider the vehicle steering system with transfer function:
\[
P(s) = \gamma\frac{s + 1/\gamma}{s^2} = \frac{s + 0.2}{s^2}
\]
where $\gamma = 5$ m is a characteristic length and the zero is at $s = -0.2$ rad/s (slow).

A state feedback controller with observer is designed to place all closed-loop poles at $s = -5$. The resulting controller transfer function is:
\[
C(s) = \frac{125s + 250}{s^2 + 20s + 100}
\]

(a) Calculate the loop transfer function $L(s)$.

(b) Determine the complementary sensitivity function $T(s)$.

(c) Find the maximum of $|T(\ii\omega)|$ and identify if there is a robustness issue.

(d) Suggest a modification to improve robustness.

\paragraph{Solution:}

\textbf{Step 1: Loop transfer function}

\[
L(s) = P(s)C(s) = \frac{s + 0.2}{s^2} \cdot \frac{125s + 250}{s^2 + 20s + 100}
\]
\[
= \frac{(s + 0.2)(125s + 250)}{s^2(s^2 + 20s + 100)} = \frac{125(s + 0.2)(s + 2)}{s^2(s^2 + 20s + 100)}
\]

\textbf{Step 2: Complementary sensitivity function}

The complementary sensitivity is:
\[
T(s) = \frac{L(s)}{1 + L(s)} = \frac{P(s)C(s)}{1 + P(s)C(s)}
\]

The closed-loop poles are all at $s = -5$ (quadruple pole), so:
\[
1 + L(s) = \frac{(s+5)^4}{s^2(s^2 + 20s + 100)}
\]

Therefore:
\[
T(s) = \frac{125(s + 0.2)(s + 2)}{(s+5)^4}
\]

\textbf{Step 3: Analyze maximum of $|T(\ii\omega)|$}

The numerator has zeros at $s = -0.2$ and $s = -2$.
The denominator has a quadruple pole at $s = -5$.

At low frequencies ($\omega \ll 0.2$), the magnitude behaves as:
\[
|T(\ii\omega)| \approx \frac{125 \cdot 0.2 \cdot 2}{5^4} = \frac{50}{625} = 0.08
\]

At the slow zero frequency $\omega = 0.2$:
\[
T(\ii \cdot 0.2) = \frac{125(\ii \cdot 0.2 + 0.2)(\ii \cdot 0.2 + 2)}{(\ii \cdot 0.2 + 5)^4}
\]

This is complex to evaluate analytically, but we can observe that the slow zero at $s = -0.2$ will cause a peak in $|T|$ near $\omega = 0.2$ rad/s.

By numerical evaluation or frequency response analysis, the maximum occurs near $\omega \approx 0.2$ rad/s with $M_t \approx 3.5$ (very large!).

\textbf{Step 4: Robustness issue and improvement}

\textbf{Problem:}
\begin{itemize}
\item The slow stable zero at $s = -0.2$ is much slower than the closed-loop poles at $s = -5$
\item This creates a large peak in the complementary sensitivity: $M_t \approx 3.5$
\item From robust stability: $|\delta| < 1/M_t \approx 0.29$, tolerating only 29\% multiplicative uncertainty
\item The system has poor robustness to high-frequency unmodeled dynamics
\end{itemize}

\textbf{Improvement:}
Instead of placing all poles at $s = -5$, place one closed-loop pole near the slow zero:

Modified pole placement: poles at $s = -0.2, -5, -5, -5$

This creates a controller pole that cancels (or nearly cancels) the process zero:
\[
C_{\text{improved}}(s) = K\frac{(s + 5)^3}{s^2 + 20s + 100} \cdot \frac{1}{s + 0.2}
\]

The cancellation results in:
\[
T_{\text{improved}}(s) \approx \frac{K(s + 2)}{(s+5)^3}
\]

This eliminates the slow zero from $T$, reducing the complementary sensitivity peak to $M_t \approx 1.2$, greatly improving robustness.

\textbf{Key insights:}
\begin{itemize}
\item Slow stable zeros should be canceled by closed-loop poles to avoid peaks in $T$
\item Cancellation improves robustness to high-frequency uncertainty
\item Design rule: assign controller poles near slow process zeros
\end{itemize}

\section{Concept 3.9.5: Nonlinear Systems}

\subsection{Example 1: Actuation Limits in a Voice Coil Servo System}

Physical actuators have fundamental limits on force, velocity, and acceleration that constrain achievable performance.

\paragraph{Problem:}
A voice coil servo system is modeled by:
\[
m\frac{\mathrm{d}^2x}{\mathrm{d}t^2} = F = k_I I
\]
where $m = 0.5$ kg is the mass, $k_I = 10$ N/A is the motor constant, and the drive has the following limits:
\begin{itemize}
\item Maximum current: $I_{\max} = 2$ A
\item Maximum voltage: $V_{\max} = 50$ V
\end{itemize}

(a) Calculate the maximum acceleration $a_{\max}$ and maximum velocity $v_{\max}$.

(b) Determine the minimum time to move the mass a distance $\ell = 0.5$ m, starting and ending at rest.

(c) Sketch the optimal velocity and acceleration profiles.

(d) Recalculate for $\ell = 2$ m and explain the difference in the solution.

\paragraph{Solution:}

\textbf{Step 1: Maximum acceleration and velocity}

Maximum acceleration is limited by the maximum current:
\[
a_{\max} = \frac{F_{\max}}{m} = \frac{k_I I_{\max}}{m} = \frac{10 \cdot 2}{0.5} = 40 \text{ m/s}^2
\]

Maximum velocity is limited by the back-EMF voltage constraint. For a voice coil, $V = k_I v$, so:
\[
v_{\max} = \frac{V_{\max}}{k_I} = \frac{50}{10} = 5 \text{ m/s}
\]

\textbf{Step 2: Determine if velocity saturation occurs}

The transition distance where velocity saturation begins is:
\[
\ell_{\text{crit}} = \frac{v_{\max}^2}{a_{\max}} = \frac{5^2}{40} = \frac{25}{40} = 0.625 \text{ m}
\]

For $\ell = 0.5$ m: Since $\ell < \ell_{\text{crit}}$, velocity saturation does \textbf{not} occur.

\textbf{Step 3: Minimum time for $\ell = 0.5$ m (no velocity saturation)}

When velocity is not saturated, the optimal strategy is:
\begin{itemize}
\item Accelerate at maximum from $t = 0$ to $t = t_1$ (reaching midpoint at $\ell/2$)
\item Decelerate at maximum from $t = t_1$ to $t = 2t_1$ (reaching $\ell$ at rest)
\end{itemize}

Using the formula for no velocity saturation:
\[
t = 2\sqrt{\frac{\ell}{a_{\max}}} = 2\sqrt{\frac{0.5}{40}} = 2\sqrt{0.0125} = 2 \cdot 0.1118 = 0.224 \text{ s}
\]

The peak velocity reached at the midpoint is:
\[
v_{\text{peak}} = a_{\max} \cdot \frac{t}{2} = 40 \cdot 0.112 = 4.48 \text{ m/s}
\]

Since $v_{\text{peak}} = 4.48 < v_{\max} = 5$ m/s, our assumption of no velocity saturation is confirmed.

\textbf{Step 4: Minimum time for $\ell = 2$ m (with velocity saturation)}

For $\ell = 2$ m: Since $\ell > \ell_{\text{crit}} = 0.625$ m, velocity \textbf{will saturate}.

The optimal strategy now has three phases:
\begin{enumerate}
\item Accelerate at $a_{\max}$ until reaching $v_{\max}$
\item Coast at constant $v_{\max}$ (zero acceleration)
\item Decelerate at $a_{\max}$ until rest
\end{enumerate}

Time to reach maximum velocity:
\[
t_1 = \frac{v_{\max}}{a_{\max}} = \frac{5}{40} = 0.125 \text{ s}
\]

Distance covered during acceleration (and deceleration):
\[
\ell_1 = \frac{v_{\max}^2}{2a_{\max}} = \frac{25}{80} = 0.3125 \text{ m}
\]

Distance at constant velocity:
\[
\ell_2 = \ell - 2\ell_1 = 2 - 2(0.3125) = 1.375 \text{ m}
\]

Time at constant velocity:
\[
t_2 = \frac{\ell_2}{v_{\max}} = \frac{1.375}{5} = 0.275 \text{ s}
\]

Total time:
\[
t_{\text{total}} = 2t_1 + t_2 = 2(0.125) + 0.275 = 0.525 \text{ s}
\]

Using the formula:
\[
t = \frac{\ell}{v_{\max}} + \frac{v_{\max}}{a_{\max}} = \frac{2}{5} + \frac{5}{40} = 0.4 + 0.125 = 0.525 \text{ s}
\]

\textbf{Step 5: Optimal profiles}

\textbf{For $\ell = 0.5$ m (no saturation):}
\begin{itemize}
\item Acceleration: $a(t) = +40$ m/s$^2$ for $0 < t < 0.112$ s, then $a(t) = -40$ m/s$^2$ for $0.112 < t < 0.224$ s
\item Velocity: Triangular profile, reaching peak of 4.48 m/s at $t = 0.112$ s
\end{itemize}

\textbf{For $\ell = 2$ m (with saturation):}
\begin{itemize}
\item Acceleration: $a(t) = +40$ m/s$^2$ for $0 < t < 0.125$ s, $a(t) = 0$ for $0.125 < t < 0.4$ s, $a(t) = -40$ m/s$^2$ for $0.4 < t < 0.525$ s
\item Velocity: Trapezoidal profile, reaching $v_{\max} = 5$ m/s and holding it for 0.275 s
\end{itemize}

\textbf{Key insights:}
\begin{itemize}
\item Actuation limits impose fundamental constraints on achievable performance
\item Short moves are acceleration-limited; long moves are velocity-limited
\item Optimal minimum-time control uses bang-bang acceleration (maximum or minimum)
\item These fundamental limits cannot be overcome by controller design alone
\end{itemize}

\subsection{Example 2: Friction-Induced Limit Cycles}

Nonlinear friction can cause oscillations even in well-designed linear control systems.

\paragraph{Problem:}
Consider a cart-pendulum balance system with state feedback. The cart has mass $M_c = 2$ kg and experiences Coulomb friction:
\[
F_{\text{friction}} = -\mu_f M_t g \operatorname{sgn}(v)
\]
where $\mu_f = 0.1$ is the coefficient of friction, $M_t = 2.5$ kg is the total mass (cart plus pendulum), $g = 9.81$ m/s$^2$, and $v$ is cart velocity.

The state feedback controller produces force $F_c$ to balance the pendulum. At equilibrium near $v = 0$, the controller attempts to apply small corrective forces.

(a) Calculate the magnitude of the friction force.

(b) Explain why friction can cause limit cycle oscillations.

(c) What is the minimum controller force needed to overcome static friction?

(d) Suggest methods to mitigate friction effects.

\paragraph{Solution:}

\textbf{Step 1: Friction force magnitude}

The Coulomb friction force magnitude is:
\[
|F_{\text{friction}}| = \mu_f M_t g = 0.1 \times 2.5 \times 9.81 = 2.45 \text{ N}
\]

This force opposes motion: $F_{\text{friction}} = -2.45 \operatorname{sgn}(v)$

\textbf{Step 2: Mechanism of limit cycles}

Coulomb friction creates a "dead zone" where small controller forces cannot produce motion:

\begin{enumerate}
\item \textbf{Near equilibrium:} The controller tries to apply a small force $|F_c| < 2.45$ N to correct a small position error
\item \textbf{Stiction:} If $|F_c| < |F_{\text{friction}}|$, the cart remains stuck ($v = 0$) despite the applied force
\item \textbf{Error accumulation:} While stuck, the integral action (if present) or accumulated state error increases the controller force
\item \textbf{Break-away:} Eventually $|F_c|$ exceeds 2.45 N, and the cart suddenly accelerates
\item \textbf{Overshoot:} Once moving, dynamic friction may be lower, causing the cart to overshoot the target
\item \textbf{Reversal:} The controller reverses force direction to correct the overshoot
\item \textbf{Cycle repeats:} The cart gets stuck again on the opposite side, creating oscillation
\end{enumerate}

\textbf{Step 3: Minimum force to overcome friction}

To initiate motion from rest:
\[
F_{\text{min}} > |F_{\text{friction}}| = 2.45 \text{ N}
\]

For the system to respond to small disturbances or reference changes, the controller must be capable of generating forces exceeding 2.45 N. However, this creates a trade-off:
\begin{itemize}
\item High controller gain $\Rightarrow$ can overcome friction but may cause violent motions and overshoot
\item Low controller gain $\Rightarrow$ smooth control but gets stuck in friction dead zone
\end{itemize}

\textbf{Step 4: Mitigation strategies}

Several approaches can reduce friction-induced oscillations:

\textbf{1. Friction compensation:}
Add a feedforward term to cancel the friction:
\[
F_{\text{total}} = F_c + \hat{F}_{\text{friction}} = F_c + 2.45 \operatorname{sgn}(v_{\text{ref}})
\]
Challenges: Requires accurate friction model and velocity measurement; sign function is discontinuous.

\textbf{2. Dead-band in control:}
If position error is below a threshold $\epsilon$, set controller output to zero:
\[
F_c = \begin{cases}
K(x_{\text{ref}} - x) & \text{if } |x_{\text{ref}} - x| > \epsilon \\
0 & \text{if } |x_{\text{ref}} - x| \leq \epsilon
\end{cases}
\]
Allows small steady-state error to avoid constant chattering.

\textbf{3. Dither signal:}
Add a small high-frequency oscillation to "shake" the system out of stiction:
\[
F_c = F_{\text{linear}} + A\sin(\omega_d t)
\]
The dither amplitude $A$ should be small enough not to disturb performance but large enough to reduce effective friction.

\textbf{4. Smooth friction model:}
Replace $\operatorname{sgn}(v)$ with a smooth approximation:
\[
F_{\text{friction}} \approx -\mu_f M_t g \tanh(\alpha v)
\]
where $\alpha$ controls the transition steepness. This makes the system more amenable to analysis and control design.

\textbf{Key insights:}
\begin{itemize}
\item Nonlinear friction effects cannot be eliminated by linear control design
\item Friction creates fundamental limits on achievable positioning accuracy
\item Practical control systems must account for actuator nonlinearities
\item Trade-offs exist between accuracy, smoothness, and robustness
\item Advanced techniques (adaptive control, sliding mode control) can better handle friction but add complexity
\end{itemize}

\end{document}
